\section*{Declaración sobre el uso de IA}

Se utilizaron herramientas de inteligencia artificial generativa en el diseño
y la curación del dataset y de las taxonomías cerradas, en el desarrollo del
código de evaluación, en la redacción y edición del manuscrito, y en la
generación de las figuras y tablas de resultados. A esos cuatro roles se
agrega un quinto, de naturaleza distinta: un modelo de lenguaje cumple un rol
\textbf{metodológico} como juez automático (Sección \ref{sec:dos-etapas}).
Ese rol no es asistencia editorial sino parte del método experimental: sus
salidas son datos del estudio, sujetos a las mismas amenazas a la validez que
cualquier otra medición. Los resultados corresponden a la ejecución real de
los modelos evaluados y a la adjudicación real del juez. Los autores
revisaron, validaron y son responsables del contenido, el análisis y las
conclusiones, en línea con marcos de buenas prácticas propuestos
recientemente para la divulgación del uso de agentes de IA en la
investigación académica \cite{torassa2025etica}.
